\section{Introduction}\label{sec:introduction}
Large language model inference is dominated by repeated linear algebra. In decoder-only transformers such as OPT \cite{opt}, every generated token passes through query, key, value, and output projections in self-attention together with expansion and contraction layers in the feed-forward network. As model width and depth increase, these linear projections dominate both the arithmetic workload and the memory footprint of inference. The dominant response in the deployment literature has therefore been dense low-precision quantization: convert weights or activations to INT8, INT4, or mixed precision, then execute the resulting tensors with highly optimized dense GEMM kernels \cite{jacob2018quantization,llmint8,zeroquant,gptq,awq,smoothquant}. That line of work is essential, but it leaves open a different systems question: whether the model can be stored and executed natively in a bit-sliced form, so that the computation itself is driven by bitwise structure rather than by reconstructing a dense low-precision matrix before every kernel launch.

This thesis addresses that question through bit-sliced vectors (BSVs). The historical term \emph{bit-sliced indexing} belongs to the database literature \cite{oneil1997variant,rinfret2001bsi}; in the present setting the stored object is a numerical vector or matrix row, so the more appropriate term is \emph{bit-sliced vector}. The objective is not merely to compress transformer weights, and it is not merely to lower arithmetic precision. The objective is to determine whether BSVs can serve as a native representation for transformer inference while preserving an execution path that operates directly on bit slices. Achieving that objective requires a full systems stack: a representation theory for fixed-point bit slicing, a CPU reference path, GPU builders for static weights and runtime activations, CUDA kernels that compute dot products directly from slices, and an evaluation framework that can measure compression, latency, and model quality at kernel, layer, and model scales.

The work began from a bitmap-based vector algebra implemented on CPU. That earlier stage established two facts. First, bit-sliced vectors preserve useful numerical structure: dot product, summation, and multiplication can be expressed directly in terms of slice-wise operations. Second, storage reduction by itself does not guarantee usable execution cost. On CPU, the bit-sliced representation often reduced memory footprint, but the dot-product path remained too slow for the vector sizes and numerical ranges that matter for transformer inference. The GPU stage of the work therefore kept the representation semantics but redesigned the execution substrate: static weights are converted once into device-side BSV structures, runtime activations are converted into BSV query slices on the GPU, and the dominant matrix products are evaluated by weighted bitwise overlap rather than by dense floating-point multiply-accumulate.

The thesis is organized around four research questions. First, can transformer linear layers be replaced by BSV-native operators while preserving end-to-end autoregressive inference on real OPT checkpoints? Second, when BSV replacement is used, which part of the pipeline dominates runtime: static construction, dynamic query construction, memory layout, or the slice-based dot kernel itself? Third, how do sensitivity and performance vary across transformer subgraphs, especially between attention and feed-forward layers? Fourth, under what conditions does the reduction in static model size translate into a measurable speedup, and why does that advantage weaken as model size grows?

The evidence collected in this work supports a more precise claim than a simple ``compression helps'' statement. Compression already translates into speedup in some regimes. Under full BSV replacement, the \texttt{opt-125m} experiments show that the bit-sliced dot path can be faster than the dense FP16 baseline, and attention-only replacement preserves a favorable latency-quality tradeoff across larger checkpoints. The main unresolved problem is scalability. As the model grows, the cost of slice-pair arithmetic, query construction, staging, and synchronization grows faster than the current kernel family can amortize. The central systems conclusion is therefore that native BSV inference is feasible, that it already yields meaningful compression and local speedup, and that the remaining research problem is to make that advantage scale to larger models through representation, layout, and kernel co-design.

\section{Background and Related Work}\label{sec:background}
\subsection{Bit-Sliced Vector Representation}
A bit-sliced vector is a fixed-point representation in which numerical values are decomposed into aligned bit planes. Let $x \in \mathbb{R}^{n}$ denote a real-valued vector, let $\alpha > 0$ denote a scale factor, and let
\begin{equation}
    \hat{x}_i = \operatorname{round}(\alpha x_i)
\end{equation}
be the corresponding fixed-point integer at position $i$. If $B$ magnitude bits are retained, then the magnitude of $\hat{x}_i$ can be written as
\begin{equation}
    |\hat{x}_i| = \sum_{b=0}^{B-1} 2^b x_i^{(b)}, \qquad x_i^{(b)} \in \{0,1\}.
\end{equation}
A signed bit-sliced vector augments these magnitude bits with a sign bit $s_i \in \{-1,1\}$ and an existence mask $e_i \in \{0,1\}$, so the decoded value is
\begin{equation}
    \hat{x}_i = e_i s_i \sum_{b=0}^{B-1} 2^b x_i^{(b)}.
\end{equation}
At the storage level, the representation consists of a stack of aligned bitmaps
\begin{equation}
    X^{(b)}[i] = x_i^{(b)}, \qquad b = 0,1,\ldots,B-1,
\end{equation}
plus auxiliary sign and existence bitmaps. The database literature originally introduced this structure for indexable attributes \cite{oneil1997variant,rinfret2001bsi}. In the present context the same structure is reinterpreted as a vector representation suitable for neural-network weights and activations.

This representation exposes a clear space--precision tradeoff. For a machine word width $W$, the raw storage complexity of an $n$-element BSV with $B$ magnitude slices is
\begin{equation}
    \mathcal{O}\!\left((B+2)\left\lceil \frac{n}{W} \right\rceil\right)
\end{equation}
words, where the additional two terms account for sign and existence metadata. The effective footprint, however, depends strongly on slice density. Some slices are sparse and compress well, while others are dense and are better stored verbatim. Hybrid bitmap schemes exploit exactly this distinction by switching storage mode according to density and expected storage cost \cite{hybridbitmap}. The same observation remains relevant for BSVs because the cost of the representation is determined not only by the nominal bit budget, but also by the occupancy pattern of the individual slices.

\subsection{Bit-Sliced Arithmetic and Computational Cost}
The practical value of a bit-sliced representation comes from the fact that arithmetic can be expressed directly on the slices. Earlier work on bit-sliced arithmetic showed that aligned slices support not only storage and filtering, but also arithmetic operators such as addition, multiplication, and dot product \cite{rinfret2001bsi}. Adopting the same row-wise convention used later for attention, let $q$ denote one query row and let $k$ denote one key row. If these rows are encoded with $S_q$ and $S_k$ magnitude slices, respectively,
\begin{equation}
    q_i = \sum_{a=0}^{S_q-1} 2^a q_i^{(a)},
    \qquad
    k_i = \sum_{b=0}^{S_k-1} 2^b k_i^{(b)},
\end{equation}
then their scalar score becomes
\begin{align}
    q k^T
    &= \sum_{i=1}^{n} q_i k_i \\
    &= \sum_{a=0}^{S_q-1} \sum_{b=0}^{S_k-1} 2^{a+b} \sum_{i=1}^{n} q_i^{(a)} k_i^{(b)}.
\end{align}
If the slice bits are packed into machine words, the inner sum can be evaluated by bitwise intersection and population count,
\begin{equation}
    q k^T = \sum_{a=0}^{S_q-1} \sum_{b=0}^{S_k-1} 2^{a+b} \, \operatorname{popcount}\!\left(Q^{(a)} \wedge K^{(b)}\right),
\end{equation}
with additional correction terms for sign handling when signed values are used. This formulation converts a dense multiply-accumulate problem into a weighted overlap-count problem. If the reduction dimension is $D$ and the packed word width is $W$, then the dominant unsigned work per score evaluation is
\begin{equation}
    \mathcal{O}\!\left(S_q S_k \frac{D}{W}\right)
\end{equation}
bitwise intersections and popcounts. That complexity exposes the central tension in the thesis: a smaller bit budget reduces storage, but the execution cost is still multiplicative in the query and key slice counts.

\subsection{Lessons from the Earlier CPU Study}
The earlier CPU study established the first empirical speed--accuracy--compression frontier for the representation. The experiments used both synthetic vectors and transformer-derived tensors. Synthetic data were used to vary vector length, numeric range, and distributional skew, while the transformer experiments used flattened query and key tensors extracted from a 12-layer BERT-based model \cite{bert,imdb}. The value of that study is not historical only. It already revealed the three quantities that continue to govern the GPU work: representation size, arithmetic accuracy, and execution cost.

For randomly generated ten-million-element vectors, the CPU results showed that BSV dot product and BSV multiplication were consistently slower than their dense baselines, and that the gap widened as the value range increased. Summation behaved more favorably and became competitive in some wider-range regimes. The memory trend moved in the opposite direction: the bit-sliced representation often consumed substantially less memory than the corresponding dense vectors. Those measurements established the key motivation for the GPU transition. The representation already had value, but CPU execution could not exploit it effectively for the dot-product path that dominates transformer inference.

The most direct precursor to the later GPU attention experiments is the BERT-based CPU dot-product study, because it isolates the same query/key interaction that later appears in the attention-score matrix. Table~\ref{tab:cpu-bert-runtime} reports the measured CPU frontier for two operating points. In the coarser operating point, the representation uses a 6-bit budget and keeps mean runtime close to the dense baseline while delivering large memory reduction. In the finer operating point, the representation expands to 9 bits, greatly reducing numerical error but increasing runtime.

\begin{table}[H]
\centering
\caption{CPU BERT dot-product results on flattened query and key tensors over 12 layers and five runs per layer.}
\label{tab:cpu-bert-runtime}
\begin{tabular}{lrrrrr}
\toprule
Bit Budget & Mean Error (\%) & Max Error (\%) & BSV Time (ms) & Dense Time (ms) & Slowdown \\
\midrule
6-bit & 0.1682 & 0.8742 & 2.178 & 2.036 & 1.070$\times$ \\
9-bit & 0.00194 & 0.00609 & 5.104 & 2.025 & 2.520$\times$ \\
\bottomrule
\end{tabular}
\end{table}

Table~\ref{tab:cpu-bert-memory} reports the corresponding storage tradeoff. The 6-bit operating point gives the stronger compression ratio, while the 9-bit operating point preserves more numerical detail at the cost of additional slices.

\begin{table}[H]
\centering
\begingroup
\singlespacing
\small
\caption{CPU BERT memory footprint and compression ratios.}
\label{tab:cpu-bert-memory}
\begin{tabular}{lrrrr}
\toprule
Bit Budget & Mean BSV Q (MB) & Mean BSV K (MB) & Q Compression & K Compression \\
\midrule
6-bit & 6.653 & 6.961 & 5.97$\times$ & 5.71$\times$ \\
9-bit & 10.984 & 10.978 & 3.62$\times$ & 3.62$\times$ \\
\bottomrule
\end{tabular}

\vspace{3pt}
\parbox{0.72\textwidth}{\footnotesize \textit{Note.} The dense query and key vectors are 39.75 MB each in both CPU settings.}
\endgroup
\end{table}

These CPU measurements already show the same structural tradeoff that later reappears on GPU. A coarser bit budget improves compression and can keep execution cost close to the dense baseline, but a finer bit budget preserves numerical fidelity at higher cost. Figure~\ref{fig:cpu-bert-runtime-error} and Figure~\ref{fig:cpu-bert-memory-compression} summarize that frontier. The figures explicitly distinguish the 6-bit and 9-bit operating points so the comparison is readable without consulting the tables.

\begin{figure}[p]
\centering
\includegraphics[width=0.95\textwidth]{cpu_bert_runtime_error.png}
\caption{CPU BERT summary for runtime and numerical error. Left: mean BSV runtime, dense runtime, and slowdown under 6-bit and 9-bit budgets. Right: mean and maximum percentage error under the same two budgets.}
\label{fig:cpu-bert-runtime-error}
\end{figure}

\begin{figure}[p]
\centering
\includegraphics[width=0.95\textwidth]{cpu_bert_memory_compression.png}
\caption{CPU BERT summary for storage and compression. The 6-bit operating point yields the stronger compression ratio, while the 9-bit operating point preserves more numerical detail.}
\label{fig:cpu-bert-memory-compression}
\end{figure}

\subsection{Related Work in Quantized LLM Inference}
The larger LLM inference literature provides strong dense low-precision baselines. Earlier work on limited-precision deep learning, integer-only inference, and accelerator co-design already showed that arithmetic width, model quality, memory traffic, and hardware mapping are tightly coupled \cite{gupta2015limitedprecision,jacob2018quantization,lee2018unpu,mishra2020survey}. That line of work established several design axes that still organize quantized inference today: post-training quantization versus quantization-aware training, weight-only versus weight-and-activation quantization, symmetric versus asymmetric scaling, per-tensor versus per-channel scaling, and static versus dynamic calibration. These choices determine not only numerical error but also whether the resulting tensors map efficiently onto dense accelerator kernels.

At LLM scale, one of the dominant numerical issues is the presence of activation outliers and heavy-tailed channels. \texttt{LLM.int8()} isolates outlier dimensions so that most of the matrix multiply can proceed in 8-bit arithmetic while a small set of channels is handled separately \cite{llmint8}. ZeroQuant combines post-training quantization with layer-wise distillation and hardware-aware integer execution \cite{zeroquant}. GPTQ formulates post-training quantization as an error-controlled weight-only problem and uses approximate second-order information to preserve quality under aggressive compression \cite{gptq}. AWQ protects salient channels identified from activation statistics \cite{awq}, while SmoothQuant redistributes activation outliers into the weights so that W8A8 execution becomes practical on modern hardware \cite{smoothquant}. These methods differ in how they calibrate and control error, but they share a structural assumption: the execution substrate remains dense.

The present work addresses a different question. The representation is required to remain bit-sliced, and the execution path is required to remain bit-sliced as well. The central problem is therefore not whether dense INT8 or INT4 inference can be made accurate and fast. The central problem is whether a non-dense representation can preserve enough arithmetic structure to support a native high-throughput execution path. This distinction is what makes the current problem a systems problem rather than simply another dense quantization baseline.

\subsection{Transformer Computation and the Attention Score Path}
Transformers provide a direct mathematical context for the BSV dot-product problem because their dominant operations are linear projections and attention-score matrix products \cite{attention,bert,opt}. Given a hidden-state matrix $X \in \mathbb{R}^{n \times d_{\mathrm{model}}}$, an attention block first forms three projections,
\begin{align}
Q &= XW_Q, \\
K &= XW_K, \\
V &= XW_V.
\end{align}
The unnormalized attention-score matrix is then
\begin{equation}
S = \frac{QK^{T}}{\sqrt{d_k}},
\end{equation}
and, in decoder-only inference, a causal mask $M$ is applied before softmax,
\begin{equation}
A = \operatorname{softmax}(S + M).
\end{equation}
The attended output is
\begin{equation}
O = AV.
\end{equation}
With multiple heads,
\begin{align}
\mathrm{head}_i &= \operatorname{Attention}(XW_i^Q, XW_i^K, XW_i^V), \\
\operatorname{MHA}(X) &= \operatorname{Concat}(\mathrm{head}_1,\ldots,\mathrm{head}_h)W^O.
\end{align}
The transformer block then applies a position-wise feed-forward sub-block,
\begin{equation}
\operatorname{FFN}(x) = W_2\,\sigma(W_1 x + b_1) + b_2,
\end{equation}
with residual and normalization layers around the attention and feed-forward components \cite{attention,bert,opt}.

These equations explain why the thesis focuses so heavily on dot products and linear layers. The score matrix $QK^T$, the $Q/K/V/O$ projections, and the FFN expansion and contraction layers are precisely the operations in which BSV replacement either succeeds or fails. They also explain why attention-only measurements are a primary result rather than a side ablation: the earlier CPU study already targeted flattened query and key tensors, so the later GPU attention experiments extend the same computational path into full-model autoregressive inference.

\section{From CPU BSV Algebra to GPU LLM Inference}\label{sec:cpu-to-gpu}
The transition from CPU bit-sliced algebra to GPU transformer inference preserves the representation semantics but changes the execution substrate completely. The CPU stage established the representational and arithmetic basis. The GPU stage specialized that basis to the operation that dominates transformer inference: high-throughput batched dot product between runtime activations and static linear weights.

At a system level, the inference path has four layers. First, values are mapped to fixed-point integers and decomposed into bit slices. Second, static weight matrices are converted once into device-resident BSV structures, while runtime activation rows are converted dynamically into batched BSV queries. Third, a CUDA kernel evaluates the weighted overlap between query slices and key slices to produce dense output rows. Fourth, a measurement layer reports kernel-only timing, layer-level timing, and end-to-end model quality on LAMBADA. Figure~\ref{fig:system-design} illustrates this organization.

\begin{figure}[t]
\centering
\begin{tikzpicture}[
  >=Latex,
  every node/.style={font=\footnotesize, align=center},
  box/.style={
    draw,
    rounded corners=2pt,
    inner sep=4pt,
    minimum height=8mm,
    line width=0.5pt
  },
  data/.style={box, fill=yellow!18},
  stage/.style={box, fill=gray!10},
  core/.style={box, fill=blue!10},
  result/.style={box, fill=green!14},
  group/.style={
    draw,
    rounded corners=3pt,
    inner sep=7pt,
    line width=0.6pt
  }
]

% -------------------------------------------------
% Top source
% -------------------------------------------------
\node[data, text width=4.6cm, minimum height=9mm] (src) at (0,0)
{Dense weights and\\activation tensors};

% -------------------------------------------------
% CPU block
% -------------------------------------------------
\node[stage, text width=3.15cm] (cpu1) at (-3.9,-2.2)
{Fixed-point mapping\\and slice construction};

\node[core, text width=3.15cm] (cpu2) at (-3.9,-3.8)
{CPU BSV arithmetic\\dot, sum, multiply};

\node[stage, text width=3.15cm] (cpu3) at (-3.9,-5.5)
{CPU transformer study\\12 layers; 6-bit\\and 9-bit budgets};

\node[group, fit=(cpu1)(cpu2)(cpu3)] (cpug) {};
\node[font=\normalsize\bfseries, above=6pt of cpug.north] (cputitle)
{CPU BSV Foundation};

% -------------------------------------------------
% Transformer block
% -------------------------------------------------
\node[stage, text width=3.15cm] (tr1) at (3.9,-2.2)
{Linear projections\\
$Q=XW_Q,\; K=XW_K$\\
$V=XW_V$};

\node[core, text width=3.15cm] (tr2) at (3.9,-3.8)
{Attention score\\
$S = QK^T/\sqrt{d_k}$};

\node[stage, text width=3.15cm] (tr3) at (3.9,-5.5)
{Masked softmax / output\\
$A=\mathrm{softmax}(S+M)$\\
$O=AV$};

\node[group, fit=(tr1)(tr2)(tr3)] (trg) {};
\node[font=\normalsize\bfseries, above=6pt of trg.north] (trtitle)
{Transformer Computation};

% -------------------------------------------------
% GPU block
% -------------------------------------------------
\node[data, text width=2.45cm] (gx) at (-2.2,-8.2)
{Runtime\\activations $X$};

\node[data, text width=2.45cm] (gw) at (2.2,-8.2)
{Static linear\\weights};

\node[stage, text width=2.65cm] (gq) at (-2.2,-9.8)
{CUDA query builder\\dynamic BSV};

\node[stage, text width=2.65cm] (gk) at (2.2,-9.8)
{CUDA key builder\\static BSV};

\node[core, text width=3.9cm] (gd) at (0,-11.3)
{BSV dot kernel\\weighted popcount};

\node[stage, text width=2.8cm] (gl) at (0,-12.7)
{BSV linear operator};

\node[stage, text width=4.15cm] (gr) at (0,-14.2)
{Layer replacement\\attention / all / mlp};

\node[group, fit=(gx)(gw)(gq)(gk)(gd)(gl)(gr)] (gpug) {};
\node[font=\normalsize\bfseries, above=6pt of gpug.north] (gputitle)
{GPU BSV Inference Stack};

% -------------------------------------------------
% Evaluation block
% -------------------------------------------------
\node[result, text width=6.6cm] (e1) at (0,-16.0)
{Benchmark modes: \texttt{kernel\_only},\\
\texttt{linear\_e2e}, \texttt{model\_e2e} on LAMBADA};

\node[result, text width=5.8cm] (e2) at (0,-17.4)
{Outputs: model size, build time,\\
dot time, Top-1, Top-5};

% -------------------------------------------------
% Arrows
% -------------------------------------------------
\draw[->] (src.south west) -- (cpu1.north);
\draw[->] (src.south east) -- (tr1.north);

\draw[->] (gx) -- (gq);
\draw[->] (gw) -- (gk);
\draw[->] (gq) -- (gd);
\draw[->] (gk) -- (gd);
\draw[->] (gd) -- (gl);
\draw[->] (gl) -- (gr);
\draw[->] (gr) -- (e1);
\draw[->] (e1) -- (e2);

% Dashed guidance arrows
\draw[dashed,->] (cpu3.south) to[out=-90,in=160] (gd.north west);
\draw[dashed,->] (tr3.south)  to[out=-90,in=20]  (gd.north east);

\end{tikzpicture}
\caption{System organization of the BSV inference pipeline. CPU-side bit-sliced construction and transformer algebra motivate a GPU execution path in which static keys and dynamic queries are built on device, evaluated by a slice-based dot kernel, and integrated into benchmarked transformer inference.}
\label{fig:system-design}
\end{figure}

\section{Methodology}\label{sec:methodology}
\subsection{Problem Formulation}
The central methodological constraint of the thesis is that weights remain stored in bit-sliced form and the dot product is computed from bit slices rather than by reconstructing a dense low-precision matrix before the kernel launch. Let $W \in \mathbb{R}^{m \times d}$ denote a dense weight matrix and let $X \in \mathbb{R}^{q \times d}$ denote a batch of activation rows. Dense inference computes
\begin{equation}
Y = XW^T.
\end{equation}
The present work replaces this with a bit-sliced formulation in which each row of $X$ is mapped to a BSV query object and each row of $W$ is mapped to a BSV key object. The output is still a dense matrix $Y$, but the intermediate computation proceeds entirely through slice-based overlap counts and fixed significance weights. The problem is therefore a coupled design problem involving representation, data layout, and parallel execution.

\subsection{Fixed-Point Mapping, Bit Budgets, and Chunking}
The floating-point tensors are first mapped to fixed-point integers. If $x_i$ is a floating-point value and $\alpha$ is a scale factor, the mapped integer is
\begin{equation}
    \hat{x}_i = \operatorname{round}(\alpha x_i).
\end{equation}
The software implementation exposes this process through decimal-shift parameters, but the thesis reports the resulting bit budgets because those are the quantities that directly determine storage and arithmetic cost. In the main GPU operating point, queries use a 7-bit budget and keys use a 6-bit budget. In the earlier CPU study, the two principal operating points correspond to 6-bit and 9-bit budgets.

The active CUDA path also applies chunk-level decomposition along the reduction dimension. If a vector is partitioned into chunks indexed by $c$, then each chunk can be represented as
\begin{equation}
    \hat{x}_{c,i} \approx \gamma_c \sum_{b=0}^{B-1} 2^b x_{c,i}^{(b)},
\end{equation}
where $\gamma_c$ is a chunk-local scale and $x_{c,i}^{(b)}$ is the $b$-th slice bit inside the chunk. This mechanism allows limited bit budgets to cover a larger dynamic range without forcing every chunk to share a single global scale. It is one of the reasons the stable 7-bit query / 6-bit key operating point remains usable on larger transformer layers.

\subsection{Static and Dynamic Work Separation}
The weight side of the pipeline is static. Each dense linear-layer weight matrix is converted once into a device-resident BSV key structure and then reused across forward passes. The activation side is dynamic. For every forward pass, each input activation row must be converted into a BSV query object before the dot-product kernel can run. This asymmetry is fundamental to the measurement methodology, because a kernel improvement is not valuable if it is offset by higher query-build cost.

This separation also clarifies the memory model. Static BSV keys determine the compressed model footprint, while dynamic BSV queries determine runtime staging overhead. The evaluation therefore reports both static model size and the timing decomposition between build cost and dot cost.

\subsection{Weighted-Overlap Dot Product}
For one query row $q$ encoded with $S_q$ magnitude slices and one key row $k$ encoded with $S_k$ magnitude slices, the unsigned component of the BSV score computation is
\begin{equation}
    q k^T = \sum_{a=0}^{S_q-1} \sum_{b=0}^{S_k-1} 2^{a+b} \, \operatorname{popcount}\!\left(Q^{(a)} \wedge K^{(b)}\right).
\end{equation}
When chunk-local scaling is present, the final result becomes a weighted sum across chunks,
\begin{equation}
    q k^T \approx \sum_{c=1}^{C} \gamma_c^{(q)} \gamma_c^{(k)}
    \sum_{a=0}^{S_q-1} \sum_{b=0}^{S_k-1} 2^{a+b}
    \, \operatorname{popcount}\!\left(Q_c^{(a)} \wedge K_c^{(b)}\right),
\end{equation}
with additional sign logic for signed slices. If the reduction dimension is $D$, the chunk width is $K_{\mathrm{chunk}}$, and the packed word width is $W$, then the dominant arithmetic work per output element can be summarized as
\begin{equation}
\begin{aligned}
T_{\mathrm{dot}}
    &= \mathcal{O}\!\left(C\,S_q\,S_k\,\frac{K_{\mathrm{chunk}}}{W}\right) \\
    &= \mathcal{O}\!\left(S_q\,S_k\,\frac{D}{W}\right).
\end{aligned}
\end{equation}
This cost model explains why the problem remains challenging even after compression. Reducing the number of stored bytes is not enough; the kernel must still evaluate every significant slice pair.

\subsection{Software Architecture and Engineering Practice}
The software stack is organized into four layers. The first layer is a CPU reference path that defines fixed-point conversion, slice construction, and reference arithmetic. The second layer is a CUDA extension that implements static key building, dynamic query building, and slice-based dot kernels. The third layer is a PyTorch integration layer that replaces dense linear operators with BSV-native operators while preserving the transformer computation graph. The fourth layer is a benchmark layer that measures kernel-only behavior, single-layer behavior, and full-model behavior with the same set of instrumentation counters.

Several engineering practices were necessary to keep the measurements defensible. Static work and dynamic work were instrumented separately so that improvements in the dot kernel could not be confused with regressions in query construction. The same benchmark harness was used across \texttt{kernel\_only}, \texttt{linear\_e2e}, and \texttt{model\_e2e} modes so that timing comparisons remained interpretable. Scope policies (\texttt{all}, \texttt{attention}, and \texttt{mlp}) were built into the replacement path rather than added as ad hoc experiments, which made layer-group sensitivity measurable without changing the rest of the software. The overall codebase is maintained as open-source software with a separate CPU bitmap library, a PyTorch/CUDA extension, and reproducible build and benchmark scripts.

\subsection{Kernel Design Rationale}
The stable kernel family treats the output matrix as a tiled computation over query rows, output columns, and 256-bit reduction chunks. Let $T_M$ denote the number of query rows handled by one cooperative thread array, let $T_N$ denote the number of output columns handled per sweep, and let $R_{\mathrm{sweep}}$ denote the number of sweeps across the output dimension performed by the same thread block. The effective CTA tile is
\begin{equation}
    M_{\mathrm{tile}} \times N_{\mathrm{tile}} = T_M \times (R_{\mathrm{sweep}} T_N).
\end{equation}
In the principal operating point,
\begin{equation}
    T_M = 32, \qquad T_N = 32, \qquad R_{\mathrm{sweep}} = 4,
\end{equation}
so the effective output tile is $32 \times 128$ over 256-bit chunks.

This geometry was selected because it balances three competing quantities. First, larger $T_M$ improves reuse of staged query slices but increases resource pressure. Second, larger $R_{\mathrm{sweep}}$ increases output coverage per CTA but also increases register use and staging cost. Third, the chunk width must remain aligned with the packed slice representation and the hardware data-movement path. In systems terms, the kernel must balance work decomposition, slice reuse, register pressure, and staging overhead. In architecture terms, it must map a weighted bitwise-overlap computation onto the CUDA execution hierarchy while preserving coalesced movement through global memory and efficient staging through shared memory \cite{nvidia_cuda_guide,nvidia_hopper_tuning,nvidia_ptx_isa}.

The empirical design process therefore evaluated geometry not as an abstract GEMM choice, but as a BSV-specific execution choice. Smaller tiles reduce query reuse and increase launch overhead. Larger sweeps increase arithmetic coverage but raise the cost of keeping intermediate accumulators live. The stable geometry is the point that remained robust across the primary OPT checkpoints under the current bit budgets. The thesis does not claim that this geometry is globally optimal. It claims that BSV kernels require geometry selection that is derived from slice-pair arithmetic and staging cost rather than copied directly from dense GEMM practice.

\subsection{Scope Policies and Layer Sensitivity}
The evaluation uses multiple scope policies because a transformer is not a homogeneous collection of linear layers. Under the present BSV regime, attention projections and feed-forward layers do not respond equally to the same bit budget. The software therefore supports three policies:
\begin{itemize}[leftmargin=2em]
\item \texttt{scope=all}: replace all eligible linear layers except the language-model head,
\item \texttt{scope=attention}: replace only attention-related projections,
\item \texttt{scope=mlp}: replace only feed-forward expansion and contraction layers.
\end{itemize}
This separation is methodologically important because it turns layer sensitivity into a measurable property of the system rather than an inferred explanation.

A broader motivation of the thesis is heterogeneous bit allocation. Let $B_q^{(\ell)}$ and $B_k^{(\ell)}$ denote the query and key budgets used in layer $\ell$. A uniform schedule can be written as
\begin{equation}
    B_q^{(\ell)} = B_q, \qquad B_k^{(\ell)} = B_k, \qquad \forall\,\ell.
\end{equation}
The measured evidence suggests instead a type-aware schedule,
\begin{equation}
    (B_q^{(\ell)}, B_k^{(\ell)}) =
    \begin{cases}
        (B_q^{\mathrm{attn}}, B_k^{\mathrm{attn}}), & \ell \in L_{\mathrm{attn}}, \\
        (B_q^{\mathrm{ffn}}, B_k^{\mathrm{ffn}}), & \ell \in L_{\mathrm{ffn}},
    \end{cases}
\end{equation}
where $L_{\mathrm{attn}}$ and $L_{\mathrm{ffn}}$ denote attention-related and feed-forward layers, respectively. The thesis does not claim that a full layer-by-layer schedule has already been optimized. It does show that the measured data reject a uniform interpretation of bit allocation across the whole transformer.

\subsection{Measurement Strategy}
The measurement layer records build time, dot time, query count, and output-element count in addition to end-to-end model metrics. This decomposition is essential for a representation whose preprocessing cost is nontrivial. A change that accelerates the dot kernel but increases query-build overhead can still regress the full model. The evaluation therefore reports kernel-only timing, layer-level timing, and model-level timing separately, together with model size and next-token accuracy.

\section{Experimental Setup}\label{sec:experimental-setup}
\subsection{Execution Environment}
The principal GPU path targets CUDA-capable devices and was tuned primarily for Hopper-class hardware. The practical execution flow consists of network initialization when required for model downloads, CUDA environment setup, Python environment activation, and local rebuild of the extension. The stable tensor-core-oriented path uses the CUDA data-movement and scheduling mechanisms exposed by the Hopper software stack \cite{nvidia_cuda_guide,nvidia_hopper_tuning,nvidia_ptx_isa}.

\subsection{Models}
The primary scaling study uses \texttt{facebook/opt-125m}, \texttt{facebook/opt-1.3b}, and \texttt{facebook/opt-6.7b}. These checkpoints provide a practical progression from a regime where compression can already translate into dot-product speedup to regimes where the current BSV execution path becomes increasingly expensive. A 30B checkpoint was used separately as a stress test to expose sensitivity by layer group, especially under aggressive full-scope replacement.

\subsection{Dataset and Metric}
End-to-end evaluation uses the LAMBADA validation split \cite{lambada}. Each input is tokenized, right-padded or truncated to a maximum sequence length of 512, and the label is taken as the last token in the sequence. The evaluator records top-1 next-token accuracy, top-5 next-token accuracy, average forward time, BSV build time, BSV dot time, and static model size.

\subsection{Principal Operating Point}
Unless stated otherwise, the principal GPU operating point uses the following settings:
\begin{itemize}[leftmargin=2em]
\item query bit budget $= 7$,
\item key bit budget $= 6$,
\item \texttt{compress\_threshold} $= 0.5$,
\item sequence length $= 512$,
\item dense baseline type $= \mathrm{fp16}$.
\end{itemize}
The software configuration retains the parameter \texttt{decimal\_places}, but the thesis reports the derived bit budgets because they are the quantities that govern both storage and arithmetic cost.

\subsection{Benchmark Modes}
The measurement harness exposes three modes.
\begin{description}[leftmargin=2em,style=nextline]
\item[\texttt{kernel\_only}] Build keys and queries once, then measure only the BSV dot kernel against dense \texttt{torch.matmul}.
\item[\texttt{linear\_e2e}] Include per-call query construction so that a single BSV linear layer can be compared against a dense linear layer.
\item[\texttt{model\_e2e}] Evaluate the full next-token pipeline on LAMBADA and report model size, latency, and accuracy together.
\end{description}
This separation is necessary because an improvement in \texttt{kernel\_only} mode does not automatically imply an improvement in \texttt{model\_e2e} mode.

\section{Results}\label{sec:results}
\subsection{Cross-Model Memory Footprint and Dot Latency}
Table~\ref{tab:crossmodel-latency} reports the main cross-model results under the fixed 7-bit query / 6-bit key operating point. All numbers correspond to the principal GPU evaluation setting with compression threshold 0.5, sequence length 512, and FP16 as the dense baseline type.

\begin{table}[H]
\centering
\caption{Cross-model static model size and per-sample dot latency under the fixed 7-bit query / 6-bit key operating point.}
\label{tab:crossmodel-latency}
\begin{tabular}{lrrrrrrr}
\toprule
Model & FP16 MB & BSV All MB & BSV Attn MB & BSV All Dot & Torch All Dot & BSV Attn Dot & Torch Attn Dot \\
\midrule
OPT-125M & 238.875 & 137.783 & 205.195 & 1.913 & 3.044 & 0.890 & 2.158 \\
OPT-1.3B & 2509.609 & 1070.453 & 2029.984 & 13.859 & 6.955 & 5.158 & 4.210 \\
OPT-6.7B & 12700.031 & 5022.281 & 10141.031 & 66.452 & 15.493 & 23.003 & 5.357 \\
\bottomrule
\end{tabular}
\end{table}

These results show that BSV compression already translates into execution benefit at some scales. Under \texttt{scope=all}, the \texttt{opt-125m} checkpoint reduces static size from 238.875 MB to 137.783 MB while also reducing per-sample dot latency from 3.044 ms to 1.913 ms, which is approximately a $1.59\times$ speedup. Under \texttt{scope=attention}, the same checkpoint reaches 0.890 ms versus 2.158 ms, or approximately a $2.42\times$ speedup on the dot path. The central scalability question is therefore not whether compression can ever reduce latency. The measured answer is yes. The systems question is why that advantage weakens as the model grows.

Figure~\ref{fig:memory-comparison} and Figure~\ref{fig:dot-latency-comparison} visualize the same trend. The memory figure decomposes the three checkpoints into separate panels so that the 125M model can be read on the same page without being visually dominated by the larger models.

\begin{figure}[H]
\centering
\includegraphics[height=0.86\textheight,keepaspectratio]{bsv_memory_comparison.png}
\caption{Static model size for OPT-125M, OPT-1.3B, and OPT-6.7B. Each checkpoint is shown in a separate panel with its own linear vertical scale so that the smallest model remains readable next to the larger checkpoints.}
\label{fig:memory-comparison}
\end{figure}

\begin{figure}[H]
\centering
\includegraphics[width=0.98\textwidth]{bsv_dot_latency_comparison.png}
\caption{Per-sample dot latency across dense Torch and BSV under \texttt{scope=all} and \texttt{scope=attention}.}
\label{fig:dot-latency-comparison}
\end{figure}

\subsection{Accuracy under Full Replacement}
Table~\ref{tab:scope-all-acc} reports top-1 and top-5 next-token accuracy when all eligible linear layers are replaced by BSV-based operators.

\begin{table}[H]
\centering
\caption{Top-1 and Top-5 accuracy under \texttt{scope=all}.}
\label{tab:scope-all-acc}
\begin{tabular}{lrrrr}
\toprule
Model & Dense Top-1 & BSV Top-1 & Dense Top-5 & BSV Top-5 \\
\midrule
OPT-125M & 0.605 & 0.626 & 0.725 & 0.729 \\
OPT-1.3B & 0.722 & 0.653 & 0.873 & 0.796 \\
OPT-6.7B & 0.798 & 0.741 & 0.929 & 0.896 \\
\bottomrule
\end{tabular}
\end{table}

The smallest model is relatively tolerant to the fixed 7-bit query / 6-bit key budget. In that regime, the reduction in model size is accompanied by dot-product speedup and little quality loss. The gap widens at 1.3B and widens further at 6.7B, indicating that the current bit budget and kernel design do not scale uniformly with model size.

\subsection{Accuracy under Attention-Only Replacement}
Table~\ref{tab:scope-attn-acc} reports the same metric when only attention-related projections are replaced.

\begin{table}[H]
\centering
\caption{Top-1 and Top-5 accuracy under \texttt{scope=attention}.}
\label{tab:scope-attn-acc}
\begin{tabular}{lrrrr}
\toprule
Model & Dense Top-1 & BSV Top-1 & Dense Top-5 & BSV Top-5 \\
\midrule
OPT-125M & 0.605 & 0.608 & 0.725 & 0.722 \\
OPT-1.3B & 0.722 & 0.666 & 0.873 & 0.821 \\
OPT-6.7B & 0.798 & 0.779 & 0.929 & 0.921 \\
\bottomrule
\end{tabular}
\end{table}

Attention-only replacement preserves quality far better than full replacement, especially at larger model sizes. This is a core empirical result of the thesis. It shows that the transformer should not be treated as a uniform collection of interchangeable linear layers. Under the present BSV regime, the attention path is substantially more tolerant than the feed-forward path.

\begin{figure}[H]
\centering
\includegraphics[width=0.98\textwidth]{bsv_accuracy_scope_all.png}
\caption{Accuracy comparison between dense FP16 and BSV under \texttt{scope=all}.}
\label{fig:accuracy-scope-all}
\end{figure}

\begin{figure}[H]
\centering
\includegraphics[width=0.98\textwidth]{bsv_accuracy_scope_attention.png}
\caption{Accuracy comparison between dense FP16 and BSV under \texttt{scope=attention}.}
\label{fig:accuracy-scope-attention}
\end{figure}

\subsection{Case Study: \texttt{opt-6.7b}}
The \texttt{opt-6.7b} checkpoint is the most informative single model in the study because it is large enough for the kernel bottleneck to be visible, but still small enough to run repeatedly during design iteration. Under the principal operating point, the dense FP16 baseline yields top-1 accuracy 0.798, top-5 accuracy 0.929, and per-sample dot latency 15.493 ms. Under full BSV replacement, the static size drops to 5022.281 MB, top-1 accuracy becomes 0.741, top-5 accuracy becomes 0.896, and per-sample dot latency rises to 66.452 ms. Under attention-only replacement, the static size is 10141.031 MB, top-1 accuracy is 0.779, top-5 accuracy is 0.921, and per-sample dot latency is 23.003 ms.

This checkpoint captures the core tradeoff of the thesis. Broader BSV replacement yields stronger compression, but it also exposes more of the model to slice-based execution cost and quantization error. Selective replacement gives up some of the compression benefit, yet preserves a much better latency-quality balance. The problem is therefore not whether BSVs can work at all. The problem is how to extend the smaller-model and attention-path advantages to larger models and broader layer scopes.

\section{Ablation Studies}\label{sec:ablation}
\subsection{Scope Ablation}
The contrast between \texttt{scope=all} and \texttt{scope=attention} is one of the strongest ablations in the thesis. It shows that the current BSV regime is not uniformly compatible with all transformer subgraphs. When all eligible layers are replaced, the model achieves the largest memory reduction but also incurs the largest latency and quality penalties. When only attention projections are replaced, the model retains much more of the dense baseline behavior while still reducing storage and dot cost relative to the dense path at small scale and relative to full-scope BSV replacement at large scale. This behavior is consistent with the broader low-bit literature, where the sensitivity of a layer depends strongly on outlier structure, activation statistics, and how much error the surrounding computation can absorb \cite{llmint8,awq,smoothquant}.

\subsection{Larger-Model Sensitivity}
A 30B checkpoint was used as a stress test to reveal how the current fixed-budget policy fails. Attention-only replacement remained relatively close to the dense baseline, while MLP-only replacement degraded more severely. When all eligible layers were replaced, the quality drop became large enough to confirm that feed-forward layers are the dominant bottleneck under the present uniform bit-budget policy.

\begin{table}[H]
\centering
\caption{Layer-group sensitivity study on \texttt{facebook/opt-30b} with 200 LAMBADA validation samples.}
\label{tab:30b-sensitivity}
\begin{tabular}{lrrrrr}
\toprule
Scope & Top-1 & Top-5 & Avg Fwd (ms) & Build (ms) & Static Size (MB) \\
\midrule
Attention & 0.820 & 0.945 & 198.878 & 27.960 & 45414.523 \\
MLP & 0.740 & 0.850 & 304.767 & 7.006 & 33655.180 \\
All & 0.470 & 0.620 & 404.781 & 29.826 & 21897.805 \\
\bottomrule
\end{tabular}
\end{table}

These results provide direct evidence for heterogeneous bit allocation. A uniform policy can compress the model aggressively, but it compresses the wrong parts too aggressively. The data suggest a hierarchy: attention layers can tolerate lower budgets earlier, while feed-forward layers require more conservative budgets if quality is to be preserved.

\subsection{Implications for Heterogeneous Bit Budgets}
A central motivation of the thesis is to move beyond one global bit budget for all layers. The measured results already indicate the structure of a more plausible policy. Let $L_{\mathrm{attn}}$ denote attention-related layers and let $L_{\mathrm{ffn}}$ denote feed-forward layers. A uniform schedule can be written as
\begin{equation}
    B_q^{(\ell)} = B_q, \qquad B_k^{(\ell)} = B_k, \qquad \forall\,\ell.
\end{equation}
The measured data suggest instead a type-aware schedule,
\begin{equation}
    (B_q^{(\ell)}, B_k^{(\ell)}) =
    \begin{cases}
        (B_q^{\mathrm{attn}}, B_k^{\mathrm{attn}}), & \ell \in L_{\mathrm{attn}}, \\
        (B_q^{\mathrm{ffn}}, B_k^{\mathrm{ffn}}), & \ell \in L_{\mathrm{ffn}}.
    \end{cases}
\end{equation}
The thesis does not claim that a full per-layer schedule has already been optimized. It does show that the current data reject the assumption that all transformer layers should share the same budget. The measured path forward is therefore to preserve higher budgets in the more sensitive feed-forward path and to lower budgets earlier in the more robust attention path, that is, $B_q^{\mathrm{ffn}} \geq B_q^{\mathrm{attn}}$ and $B_k^{\mathrm{ffn}} \geq B_k^{\mathrm{attn}}$.

\subsection{Lower-Bit Frontier}
The lower-bit frontier clarifies what it means to accelerate the system by reducing the slice count. On \texttt{opt-6.7b}, moving from a 7-bit query / 6-bit key operating point to a 6-bit query / 5-bit key operating point improves both latency and storage, but it also produces a clear quality drop.

\begin{table}[H]
\centering
\caption{Lower-bit frontier study on \texttt{facebook/opt-6.7b} using 200 LAMBADA validation samples.}
\label{tab:67b-frontier}
\begin{tabular}{lrrrrr}
\toprule
Configuration & Top-1 & Top-5 & Dot (ms) & Avg Fwd (ms) & Static Size (MB) \\
\midrule
7-bit query / 6-bit key & 0.745 & 0.920 & 67.952 & 123.237 & 5049.281 \\
6-bit query / 5-bit key & 0.680 & 0.840 & 54.226 & 109.031 & 4276.781 \\
\bottomrule
\end{tabular}
\end{table}

The latency reduction is real, but it is purchased by moving to a different numerical operating point. This is why the thesis does not interpret lower-bit speedup as a kernel-design victory at fixed quality. The more meaningful systems objective is to improve execution efficiency at a fixed or comparable operating point, while treating lower-bit settings as a separate speed--accuracy frontier.

\subsection{CPU-to-GPU Continuity}
One of the useful ablations in the thesis is not a parameter sweep but a continuity argument across execution substrates. The CPU study already showed that the representation could be compact and numerically useful, yet its dot-product path was too slow for large neural workloads. The GPU study does not replace that conclusion; it extends it. At small scale, the same representation can now yield real speedup on the dot path. At larger scale, the same representation still incurs more slice interaction and staging cost than the current kernels can absorb. This continuity matters because it shows that the thesis is not comparing unrelated systems. It is tracing the same representation across two execution substrates and identifying where the performance bottleneck migrates.

\section{Discussion}\label{sec:discussion}
The thesis establishes a complete path from bit-sliced vector representation to end-to-end transformer inference. The CPU stage demonstrates that the representation is algebraically meaningful and storage-efficient. The GPU stage demonstrates that the representation can be integrated into a real autoregressive inference pipeline and measured at kernel, layer, and model scales. That progression matters because it turns the problem from a data-structure exercise into a systems problem.

The measurements support two conclusions simultaneously. First, BSV compression already yields real runtime benefit in some regimes. The 125M checkpoint shows faster dot execution under full BSV replacement, and attention-only replacement preserves a useful latency-quality balance across larger models. Second, that advantage does not yet scale cleanly with model size. The dot path becomes progressively more expensive at 1.3B and 6.7B scale, and the gap is most severe when the same fixed budget is applied indiscriminately to all linear layers.

The attention-only and 30B sensitivity results sharpen that conclusion further. They show that the problem is not simply ``bit-sliced arithmetic is too slow.'' Instead, the behavior depends strongly on which subgraph is replaced. Attention layers tolerate the current BSV regime substantially better than feed-forward layers, which means that the representation, the layer type, and the execution path must be co-designed. In other words, the remaining problem is not feasibility but scalability.

This framing is what makes the systems contribution defensible. The thesis does not present BSVs as a finished replacement for dense FP16 inference. It presents a measured operating region in which compression already turns into speedup, a larger operating region in which compression still works but speedup does not scale, and a set of empirical observations that explain why. That is enough to support a concrete research direction rather than a vague intuition.

\section{Future Work}\label{sec:future-work}
The next stage of the work follows directly from the empirical structure already observed. The most immediate direction is heterogeneous bit allocation. The data already show that attention and feed-forward layers should not share the same budget. The natural next step is therefore a type-aware or layer-aware schedule that preserves more slices in the most sensitive layers and reduces them more aggressively where the model is more tolerant.

A second direction is improved kernel families at fixed quality. The current kernels still evaluate many slice-pair interactions per chunk, and the cost of that arithmetic grows rapidly with model size. Future work should therefore target arithmetic formulations and data layouts that reduce effective slice interaction cost without abandoning the native BSV execution model.

A third direction is broader model coverage. OPT and LAMBADA provide a clean baseline for the present study, but the same methodology should be applied to additional transformer families and broader evaluation suites. This would clarify whether the observed attention-versus-MLP sensitivity is specific to OPT or reflects a more general property of BSV inference.

\section{Conclusion}\label{sec:conclusion}
This work establishes a complete path from bit-sliced vector representation to GPU-based large language model inference. The earlier CPU stage showed that bit-sliced vectors can represent numerical data compactly and support arithmetic directly in slice form, but that CPU execution was insufficient for the dot-product path that dominates transformers. The GPU stage preserved the same bit-sliced semantics and scaled them into a full inference stack: static weights were converted into BSV key structures, runtime activations were converted into BSV query slices, and transformer linear layers were executed through fixed-budget CUDA dot kernels.

The main conclusion is concrete. Compression already yields measurable speedup in some regimes, especially at smaller model scale and under selective replacement of attention layers. The unresolved challenge is to make that advantage scale to larger models. The thesis therefore identifies not only a working BSV-native inference path, but also the precise systems problem that remains: how to co-design bit budget, layout, and kernel structure so that the storage advantage of BSVs continues to translate into execution advantage as transformer scale increases.
